\section{Methodology}
\label{sec:method}

\subsection{Training and evaluation pipeline}
\reporttodo{Add a diagram of the form raw paired data, preprocessing, data
loader, encoders, fusion, loss, optimization, prediction, evaluation (Handbook
Section 31). Explain each stage in one or two sentences and name the main source
files.}

\reportfigure{Images/a3_pipeline.png}{Training and evaluation pipeline for
image-text retrieval on Flickr8k.}{fig:pipeline}

\subsection{Modality encoders and representations}
\label{sec:encoders}
\reporttodo{For each modality, describe the encoder, the representation it
produces (for example, a pooled embedding of dimension $d$), the pretrained
weights and library, and which layers are trained. Explain how an image and a
caption are aligned: every pair comes from the same Flickr8k image and each
image has five captions.}

\subsection{Unimodal baselines}
\label{sec:unimodal}
\reporttodo{Describe the baseline that uses only the image modality and the
baseline that uses only the text modality (Handbook Section 29.1). State
exactly how each one produces a ranking for retrieval, which information it
cannot use, and what its result shows. List the strengths and limits of each
baseline for this task. Give parameter counts and state which parts were
written by the group and which come from libraries (name the libraries).}

\subsection{Simple fusion baseline}
\label{sec:simple-fusion}
\reporttodo{Describe the simple fusion model: the fusion stage (early,
intermediate or late), the fusion operation, and how it scores an image-caption
pair. Explain why it is the simplest reasonable way to combine the two
modalities, and list its strengths and limits for this task. Give its parameter
count and state which parts were written by the group and which come from
libraries.}

\subsection{Improved multimodal model}
\label{sec:improved}
\reporttodo{Describe the improved model or fusion design: what changes from the
simple fusion baseline, why the change should help retrieval, and its pretraining
and fine-tuning strategy. List its strengths and limits for this task. Give its
parameter count, the source of each pretrained component, and which parts were
written by the group.}

\begin{table}[htbp]
\centering\small
\caption{Compared models. Implementation states whether the group wrote the
model or used a library.}
\label{tab:model-summary}
\begin{tabularx}{\linewidth}{lYllr}
\toprule
Model & Role & Initialization & Implementation & Parameters\\
\midrule
\pending & Image-only baseline & \pending & \pending & \pending\\
\pending & Text-only baseline & \pending & \pending & \pending\\
\pending & Simple fusion baseline & \pending & \pending & \pending\\
\pending & Improved multimodal model & \pending & \pending & \pending\\
\bottomrule
\end{tabularx}
\end{table}


\subsection{Loss and evaluation metrics}
\label{sec:metrics}
\reporttodo{Define the per-modality losses, if any, and the overall training
loss, for example a symmetric contrastive loss over a batch of $B$ pairs with
temperature $\tau$. Define Recall@$K$ for $K\in\{1,5,10\}$ and the median rank
for both directions. State that image-to-text retrieval counts a hit when any of
the image's five captions is in the top $K$, and that each caption has exactly
one correct image. Name the validation metric used for checkpoint selection.}

\subsection{Missing-modality handling}
\label{sec:missing}
\reporttodo{State that every Flickr8k image has five captions and every caption
has its image, so no sample lacks a modality. Say whether any model is tested
with one modality removed, and how.}

\subsection{Data loading and training loop}
\label{sec:loop}
\reporttodo{Describe how the data loader builds a batch of image-caption pairs:
how one of the five captions is chosen for each image, how images and captions
stay aligned, and how padding is applied. Describe one training step: encoding
both modalities, fusion, the loss over the batch (for example, with the other
pairs in the batch as negatives), and the optimizer update. Describe how
validation encodes all images and captions once and ranks them. Name the source
files.}

\subsection{Training configuration}
\label{sec:training}
\reporttodo{Fill Table~\ref{tab:training-setup} from the configs in
\codepath{configs/} and the saved launch configs. Explain how hyperparameters
were chosen. Note any difference between a config file and the reported run.}

\begin{table}[htbp]
\centering\small
\caption{Training settings for the reported runs. Add one column per trained
model.}
\label{tab:training-setup}
\begin{tabular}{lll}
\toprule
Setting & \pending & \pending\\
\midrule
Image input size & \pending & \pending\\
Text tokenizer / max length & \pending & \pending\\
Trainable layers & \pending & \pending\\
Optimizer & \pending & \pending\\
Learning rate & \pending & \pending\\
Scheduler & \pending & \pending\\
Batch size (effective) & \pending & \pending\\
Epochs / seed & \pending & \pending\\
Weight decay & \pending & \pending\\
Other regularization & \pending & \pending\\
Loss / temperature & \pending & \pending\\
Mixed precision & \pending & \pending\\
Early stopping & \pending & \pending\\
Checkpoint rule & \pending & \pending\\
Hardware & \pending & \pending\\
Software versions & \pending & \pending\\
\bottomrule
\end{tabular}
\end{table}


\subsection{Experimental design}
\label{sec:experiments}
\reporttodo{For each experiment, state the hypothesis, the changed factor, the
fixed factors, the metrics, and the decision criterion (Handbook Section 30).
The comparisons are the two unimodal baselines against the simple fusion
baseline, the simple fusion baseline against the improved model, and at least
one ablation that removes or changes a modality or a fusion component. Define in
advance what counts as a meaningful difference given the number of seeds.}

\begin{table}[htbp]
\centering\small
\caption{Experimental design. Each row changes one factor and fixes the rest.}
\label{tab:experiment-design}
\begin{tabularx}{\linewidth}{lYYYYY}
\toprule
Experiment & Hypothesis & Changed factor & Fixed factors & Metrics & Decision criterion\\
\midrule
Unimodal vs.\ fusion & \pending & \pending & \pending & \pending & \pending\\
Simple vs.\ improved fusion & \pending & \pending & \pending & \pending & \pending\\
Ablation & \pending & \pending & \pending & \pending & \pending\\
\bottomrule
\end{tabularx}
\end{table}

